\hypertarget{advanced-text-vs-bytes-unicode-internals}{%
\chapter{ADVANCED TEXT VS BYTES \& UNICODE
INTERNALS}\label{advanced-text-vs-bytes-unicode-internals}}

\hypertarget{the-pre-pep-393-memory-overhead}{%
\subsection{8.1 The Pre-PEP 393 Memory
Overhead}\label{the-pre-pep-393-memory-overhead}}

Historically, Python 3 represented all Unicode strings using either
\textbf{UCS-2} (2 bytes per character) or \textbf{UCS-4} (4 bytes per
character) wide character arrays, configured at compilation time. *
\textbf{UCS-2 Pitfall}: Limited to characters under 65535. Characters
beyond that required surrogate pairs, making indexing (\(O(1)\) lookup
by offset) complex. * \textbf{UCS-4 Pitfall}: Consumed 4 bytes per
character regardless of contents. An ASCII string like
\passthrough{\lstinline!"hello"!} (which requires only 5 bytes) consumed
20 bytes on the heap, wasting memory.

\hypertarget{pep-393-flexible-string-representation}{%
\subsection{8.2 PEP 393: Flexible String
Representation}\label{pep-393-flexible-string-representation}}

To solve memory overhead, Python 3.3 introduced \textbf{PEP 393}.
CPython now selects the most compact memory layout based on the string's
characters:

\begin{lstlisting}
PEP 393 String Header and Struct Formats:
1. PyASCIIObject:
   +-------------------------------------------------------------+
   | ob_refcnt | ob_type | length | state (ASCII=1, compact=1)  | -> Raw ASCII data (1 byte/char)
   +-------------------------------------------------------------+

2. Compact 1-Byte (Latin-1):
   +-------------------------------------------------------------+
   | ob_refcnt | ob_type | length | state (ASCII=0, compact=1)  | -> Latin-1 data (1 byte/char)
   +-------------------------------------------------------------+

3. Compact 2-Byte (UCS-2):
   +-------------------------------------------------------------+
   | ob_refcnt | ob_type | length | state (UCS-2)               | -> UCS-2 data (2 bytes/char)
   +-------------------------------------------------------------+

4. Compact 4-Byte (UCS-4):
   +-------------------------------------------------------------+
   | ob_refcnt | ob_type | length | state (UCS-4)               | -> UCS-4 data (4 bytes/char)
   +-------------------------------------------------------------+
\end{lstlisting}

\begin{itemize}
\tightlist
\item
  \textbf{Runtime Memory Analysis}: We can verify this dynamic memory
  layout using \passthrough{\lstinline!sys.getsizeof()!}:
\end{itemize}

\begin{lstlisting}[language=Python]
import sys

# ASCII String: 1 byte per character + 49 bytes header
str_ascii = "hello"
print("ASCII size:", sys.getsizeof(str_ascii), "bytes") # 49 + 5 = 54

# Latin-1 String (Unicode character <= 255): 1 byte per character + 73 bytes header (non-ASCII compact)
str_latin = "hell"
print("Latin-1 size:", sys.getsizeof(str_latin), "bytes") # 73 + 5 = 78

# UCS-2 String (Unicode character <= 65535): 2 bytes per character + 73 bytes header
str_ucs2 = "hell"
print("UCS-2 size:", sys.getsizeof(str_ucs2), "bytes") # 73 + (6 * 2) = 85

# UCS-4 String (Unicode character > 65535, e.g. emoji): 4 bytes per character + 73 bytes header
str_ucs4 = "hell"
print("UCS-4 size:", sys.getsizeof(str_ucs4), "bytes") # 73 + (7 * 4) = 101
\end{lstlisting}

\hypertarget{encoding-decoding-and-binary-buffers}{%
\subsection{8.3 Encoding, Decoding, and Binary
Buffers}\label{encoding-decoding-and-binary-buffers}}

\begin{itemize}
\tightlist
\item
  \textbf{The Boundary}: Python strictly enforces the boundary between
  text (\passthrough{\lstinline!str!}, Unicode code points) and binary
  data (\passthrough{\lstinline!bytes!}, raw 8-bit octets).

  \begin{itemize}
  \tightlist
  \item
    \passthrough{\lstinline!str.encode(encoding)!} translates Unicode
    code points into raw byte sequences.
  \item
    \passthrough{\lstinline!bytes.decode(encoding)!} decodes raw bytes
    back into Unicode strings.
  \end{itemize}
\item
  \textbf{Zero-Copy Slicing via memoryview}:
  \passthrough{\lstinline!memoryview!} objects allow sharing access to
  binary buffers (like \passthrough{\lstinline!bytearray!} or file
  buffers) without copying the underlying bytes. This is critical for
  high-performance networking and I/O tasks.
\end{itemize}

\begin{lstlisting}[language=Python]
# Zero-copy slicing demo
data = bytearray(b"system_payload_data")
view = memoryview(data)

# Slice without copying
sub_view = view[7:14]
print("Slice contents:", sub_view.tobytes())

# Modify buffer in-place
sub_view[0] = ord(b"P")
print("Modified original array:", data) # data becomes bytearray(b"system_Payload_data")
\end{lstlisting}

\begin{center}\rule{0.5\linewidth}{0.5pt}\end{center}
